\documentclass[../../../include/open-logic-section]{subfiles}

\begin{document}

\olfileid[id]{sth}{story}{predicative}	

\olsection{Predikatif dan Impredikatif}

Himpunan Russell, $R$, didefinisikan melalui $\Setabs{x}{x \notin
x}$. Jika diuraikan lebih lengkap, $R$ adalah himpunan yang memuat semua dan
hanya himpunan yang bukan anggota dirinya sendiri. Jadi, ketika mendefinisikan
$R$, kita menguantifikasi atas domain yang akan memuat $R$ (seandainya objek
tersebut ada).

Ini merupakan definisi \emph{impredikatif}. Secara lebih umum, kita dapat
mengatakan bahwa suatu definisi impredikatif jika dan hanya jika definisi itu menguantifikasi
atas domain yang memuat objek yang sedang didefinisikan.
	
Sesudah munculnya berbagai paradoks, Whitehead, Russell, Poincar\'{e}, dan
Weyl menolak definisi impredikatif semacam itu karena membentuk
``lingkaran setan'':
\begin{quote}
	Analisis atas paradoks-paradoks yang harus dihindari menunjukkan bahwa
	semuanya muncul dari sejenis lingkaran setan. Lingkaran-lingkaran setan
	tersebut timbul dari anggapan bahwa suatu koleksi objek dapat
	memuat anggota yang hanya dapat didefinisikan melalui koleksi itu secara
	keseluruhan[\ldots. \textparagraph]

	Prinsip yang memungkinkan kita menghindari totalitas yang tidak sah dapat
	dinyatakan sebagai berikut: `Apa pun yang melibatkan \emph{semua} anggota
	suatu koleksi tidak boleh menjadi salah satu anggota koleksi itu'; atau,
	sebaliknya: `Jika, seandainya suatu koleksi tertentu mempunyai suatu
	totalitas, koleksi itu akan mempunyai anggota yang hanya dapat didefinisikan
	dengan merujuk pada totalitas tersebut, maka koleksi itu tidak mempunyai
	totalitas.' Kita akan menyebut ini `prinsip lingkaran setan', karena
	prinsip ini memungkinkan kita menghindari lingkaran-lingkaran setan yang
	terlibat dalam pengandaian totalitas yang tidak sah.
	\citep[p.~37]{WhiteheadRussell1910}
\end{quote}
Jika kita mengikuti mereka dengan menerima \emph{prinsip lingkaran setan},
kita dapat mencoba mengganti Skema Komprehensi Naif yang membawa bencana
(dari \olref[sth][story][rus]{sec}) dengan sesuatu seperti berikut:

\begin{defish}
\emph{Komprehensi Predikatif.} Untuk setiap formula $\phi$ yang hanya
menguantifikasi atas himpunan: himpunan$^\prime$
$\Setabs{x}{\phi(x)}$ ada.
\end{defish}

Selama himpunan$^{\prime}$ bukan merupakan himpunan, tidak akan timbul
kontradiksi.

Sayangnya, Komprehensi Predikatif tidak terlalu \emph{komprehensif}. Bagaimanapun,
prinsip itu memperkenalkan entitas baru kepada kita, yakni
himpunan$^\prime$. Jadi, kita harus mempertimbangkan formula yang
menguantifikasi atas himpunan$^\prime$. Jika formula semacam itu selalu
menghasilkan himpunan$^\prime$, Paradoks Russell akan timbul lagi hanya dengan
mempertimbangkan himpunan$^\prime$ dari semua himpunan$^\prime$ yang bukan
anggota dirinya sendiri. Dengan melanjutkan gagasan yang sama, kita harus
mengatakan bahwa formula yang menguantifikasi atas himpunan$^\prime$
menghasilkan himpunan$^{\prime\prime}$ yang bersesuaian. Lalu kita akan
memerlukan himpunan$^{\prime\prime\prime}$,
himpunan$^{\prime\prime\prime\prime}$, dan seterusnya. Untuk mencegah ledakan
tanda prima, lebih mudah menganggap semuanya sebagai himpunan$_0$,
himpunan$_1$, himpunan$_2$, himpunan$_3$, himpunan$_4$,\ldots. Dengan cara
ini kita memperoleh jalan menuju teori tipe (sederhana).

Ada beberapa keberatan yang jelas terhadap teori semacam itu (meskipun tidak
jelas bahwa keberatan tersebut \emph{sangat menentukan}). Singkatnya: teori
yang dihasilkan rumit digunakan; teori itu boros dalam mendalilkan beragam
jenis objek; dan pada akhirnya tidak jelas bahwa definisi impredikatif bahkan
\emph{seburuk itu}.
	
Untuk menegaskan butir terakhir, perhatikan komentar
\citeauthor{Ramsey1925} berikut:
\begin{quote}
	kita dapat menyebut seorang pria sebagai yang tertinggi dalam suatu
	kelompok, sehingga mengidentifikasinya melalui suatu totalitas yang ia
	sendiri menjadi anggotanya, tanpa adanya lingkaran setan.
	\citep{Ramsey1925}
\end{quote}
Butir yang diajukan Ramsey ialah bahwa ``pria tertinggi dalam kelompok''
\emph{merupakan} definisi impredikatif; tetapi definisi itu jelas sepenuhnya
wajar.

Seseorang mungkin menanggapi bahwa, dalam kasus ini, kita dapat memilih orang
tertinggi itu dengan cara \emph{predikatif}. Misalnya, mungkin kita cukup
menunjuk pria yang dimaksud. Karena itu, keberatan terhadap definisi
impredikatif jelas harus dibatasi pada entitas yang \emph{hanya} dapat dipilih
secara impredikatif. Namun, sekalipun demikian, kita masih memerlukan
penjelasan lebih lanjut tentang mengapa ``keimpredikatifan esensial'' semacam
itu begitu buruk.\footnote{Untuk pembahasan lebih lanjut, lihat
\citet{Linnebo2010}.}

Memang, definisi impredikatif membawa masalah yang sangat serius jika kita
ingin definisi kita menyediakan semacam resep untuk \emph{menciptakan} suatu
objek. Sebab, dengan diberikannya definisi impredikatif, kita benar-benar akan
terperangkap dalam lingkaran setan: untuk menciptakan objek yang
ditentukan secara impredikatif, kita harus \emph{terlebih dahulu} menciptakan
semua objek (termasuk objek yang ditentukan secara impredikatif), karena objek
tersebut ditentukan dengan mengacu pada semua objek; jadi, kita harus
menciptakan objek yang ditentukan secara impredikatif sebelum kita selesai
menciptakan objek itu sendiri. Namun sekali lagi, ini hanya merupakan
keberatan serius terhadap himpunan yang ditentukan ``secara impredikatif
esensial'' jika kita menganggap himpunan sebagai benda yang kita
\emph{ciptakan}. Dan (barangkali) kita tidak menganggapnya demikian.

Karena itu---baik atau buruk---pendekatan yang kemudian menjadi umum tidak
mengambil sikap keras mengenai (im)\-pre\-di\-ka\-tiv\-itas. Sebaliknya,
pendekatan tersebut menggunakan apa yang kini dianggap sebagai pendekatan
kumulatif-iteratif. Pada akhirnya, pendekatan ini memungkinkan kita
menstratifikasi himpunan ke dalam ``tahap-tahap''---\emph{agak} menyerupai
cara pendekatan predikatif menstratifikasi entitas menjadi himpunan$_0$,
himpunan$_1$, himpunan$_2$, \ldots---tetapi kita tidak akan mendalilkan
perbedaan jenis apa pun di antara semuanya.

\end{document}
